\documentclass[11pt,a4paper]{article}
\usepackage[utf8]{inputenc}
\usepackage{amsmath,amssymb}
\usepackage{geometry}
\geometry{margin=2.5cm}
\usepackage{listings}
\usepackage{xcolor}
\usepackage{hyperref}

\lstset{
  language=C++,
  basicstyle=\ttfamily\small,
  keywordstyle=\color{blue}\bfseries,
  commentstyle=\color{green!60!black},
  stringstyle=\color{red!70!black},
  numbers=left,
  numberstyle=\tiny\color{gray},
  numbersep=8pt,
  frame=single,
  breaklines=true,
  tabsize=4,
  showstringspaces=false
}

\title{IOI 1995 -- Shopping Offers}
\author{Solution}
\date{}

\begin{document}
\maketitle

\section{Problem Statement}

A shop sells $b$ different products ($b \le 5$). Each product $i$ has a unit
price $c_i$, and we wish to buy exactly $q_i$ units of product $i$ ($q_i \le 5$).

There are $s$ special offers ($s \le 99$). Each offer specifies a bundle: for each product,
how many units are included, along with a bundle price. An offer may be used as many times
as desired, but we cannot buy more than the required quantity of any product.

Find the minimum total cost to buy exactly the required quantities.

\section{Solution Approach}

\subsection{State-Space DP}

Since there are at most 5 products each with quantity at most 5, the state space is
at most $6^5 = 7{,}776$ states. Define:
\[
  \mathrm{dp}[q_1][q_2][q_3][q_4][q_5] = \text{minimum cost to buy exactly } (q_1, q_2, \ldots, q_5) \text{ items.}
\]

\subsection{Base Case}
$\mathrm{dp}[0][0][0][0][0] = 0$.

\subsection{Transitions}
For each state $(q_1, \ldots, q_b)$ with at least one $q_i > 0$:
\begin{enumerate}
  \item \textbf{Buy one unit of product $i$:} If $q_i > 0$, transition from
    $(q_1, \ldots, q_i - 1, \ldots, q_b)$ with cost increase $c_i$.
  \item \textbf{Use offer $j$:} If $q_i \ge o_{j,i}$ for all $i$, transition from
    $(q_1 - o_{j,1}, \ldots, q_b - o_{j,b})$ with cost increase equal to the offer price.
\end{enumerate}

We compute the DP bottom-up, iterating through all quantity tuples in lexicographic order.
Each state takes its value as the minimum over all valid predecessors.

\section{C++ Solution}

\begin{lstlisting}
#include <cstdio>
#include <cstring>
#include <algorithm>
using namespace std;

int b; // number of products (1..5)
int price[6]; // unit price for each product (1-indexed)
int want[6];  // desired quantity for each product
int code[6];  // product code mapping

int s; // number of offers
struct Offer {
    int qty[6]; // quantity of each product in offer (1-indexed)
    int cost;
};
Offer offers[100];

int dp[6][6][6][6][6];

int main() {
    scanf("%d", &s);

    // Read offers (product codes used directly for now)
    for (int i = 0; i < s; i++) {
        int nk;
        scanf("%d", &nk);
        memset(offers[i].qty, 0, sizeof(offers[i].qty));
        for (int j = 0; j < nk; j++) {
            int c, q;
            scanf("%d %d", &c, &q);
            offers[i].qty[c] = q; // store by product code
        }
        scanf("%d", &offers[i].cost);
    }

    scanf("%d", &b);
    for (int i = 1; i <= b; i++)
        scanf("%d %d %d", &code[i], &want[i], &price[i]);

    // Remap offer quantities from product codes to indices 1..b
    Offer mapped[100];
    for (int i = 0; i < s; i++) {
        mapped[i].cost = offers[i].cost;
        memset(mapped[i].qty, 0, sizeof(mapped[i].qty));
        for (int j = 1; j <= b; j++)
            mapped[i].qty[j] = offers[i].qty[code[j]];
    }

    // Initialize DP to infinity
    memset(dp, 0x3f, sizeof(dp));
    dp[0][0][0][0][0] = 0;

    // Set up quantity bounds
    int w[6];
    for (int i = 1; i <= 5; i++)
        w[i] = (i <= b) ? want[i] : 0;

    // Bottom-up DP over all quantity tuples
    for (int q1 = 0; q1 <= w[1]; q1++)
    for (int q2 = 0; q2 <= w[2]; q2++)
    for (int q3 = 0; q3 <= w[3]; q3++)
    for (int q4 = 0; q4 <= w[4]; q4++)
    for (int q5 = 0; q5 <= w[5]; q5++) {
        if (q1 == 0 && q2 == 0 && q3 == 0 && q4 == 0 && q5 == 0)
            continue; // base case already set

        int best = 0x3f3f3f3f;
        int qq[6] = {0, q1, q2, q3, q4, q5};

        // Option 1: buy a single unit of product j
        for (int j = 1; j <= b; j++) {
            if (qq[j] > 0) {
                int prev[6] = {0, q1, q2, q3, q4, q5};
                prev[j]--;
                int val = dp[prev[1]][prev[2]][prev[3]][prev[4]][prev[5]];
                if (val < 0x3f3f3f3f)
                    best = min(best, val + price[j]);
            }
        }

        // Option 2: use an offer
        for (int o = 0; o < s; o++) {
            bool ok = true;
            int prev[6] = {0, q1, q2, q3, q4, q5};
            for (int j = 1; j <= b; j++) {
                prev[j] -= mapped[o].qty[j];
                if (prev[j] < 0) { ok = false; break; }
            }
            if (ok) {
                int val = dp[prev[1]][prev[2]][prev[3]][prev[4]][prev[5]];
                if (val < 0x3f3f3f3f)
                    best = min(best, val + mapped[o].cost);
            }
        }

        dp[q1][q2][q3][q4][q5] = best;
    }

    printf("%d\n", dp[w[1]][w[2]][w[3]][w[4]][w[5]]);
    return 0;
}
\end{lstlisting}

\section{Correctness}

The DP processes states in non-decreasing order of the quantity tuple (lexicographically).
Every predecessor state $(q_1, \ldots, q_i - 1, \ldots, q_b)$ or
$(q_1 - o_{j,1}, \ldots, q_b - o_{j,b})$ is componentwise $\le$ the current state,
so it has already been computed when needed. The base case $\mathrm{dp}[\mathbf{0}] = 0$
is correct, and each transition adds exactly the cost of one purchase, so by induction
the DP computes the minimum cost for every reachable state.

\section{Complexity Analysis}

\begin{itemize}
  \item \textbf{Time complexity:} $O(Q^b \cdot (b + s))$, where $Q = \max(q_i) + 1 \le 6$,
        $b \le 5$, and $s \le 99$. With these bounds: $6^5 \times 104 \approx 810{,}000$ operations.
  \item \textbf{Space complexity:} $O(Q^5) = O(7{,}776)$ for the DP table.
\end{itemize}

\end{document}
